\subsection{Entregables [Pendientes: 1. acordar fechas y cantidades, 2. precisar destinatarios y responsables de cada entrega]}\label{sec:spmp-6-4}
El catálogo siguiente conserva la organización de productos empleada en el VFP y la relaciona con las evidencias previstas en la adaptación de estándares. Los identificadores permiten que las secciones posteriores indiquen criterios, actividades y versiones sin repetir la descripción del producto. La entrega académica tendrá como destinatario al director, mientras que los documentos de trabajo también servirán al equipo y las evaluaciones involucrarán a los usuarios que se acuerden. No se han confirmado cantidades de escenarios ni fechas particulares de entrega, por lo que esos valores permanecerán asociados a la planificación pendiente de 2027-I~\cite{vfp2026,adaptacion2026}.
\begin{longtable}{|p{0.09\dimexpr\linewidth-34pt\relax}|p{0.59\dimexpr\linewidth-34pt\relax}|p{0.32\dimexpr\linewidth-34pt\relax}|}
\caption{Catálogo de entregables del proyecto. Fuente: adaptación del VFP y de la metodología del equipo~\cite{vfp2026,adaptacion2026}.}\\
\hline
\textbf{ID} & \textbf{Producto y destinatario} & \textbf{Momento previsto}\\\hline
\endfirsthead
\hline
\textbf{ID} & \textbf{Producto y destinatario} & \textbf{Momento previsto}\\\hline
\endhead
E1 & Especificación viva de requisitos, alcance y trazabilidad. Equipo y director. & Preparación de Construcción y refinamiento durante el desarrollo.\\\hline
E2 & Arquitectura y modelos UML que describen responsabilidades e intercambios. Equipo y director. & Diseño inicial y actualización con cada incremento pertinente.\\\hline
E3 & Paquete reproducible de configuración WFDS, entradas preparadas e informe del piloto. Equipo y director. & Piloto y comprobaciones técnicas de Construcción.\\\hline
E4 & Biblioteca y catálogo que relacionan escenarios con condiciones, procedencia y límites. Aplicación, equipo y director. & Generación e integración durante Construcción.\\\hline
E5 & Código fuente y versión ejecutable de la aplicación de escritorio. Equipo, director y usuarios evaluadores. & Incrementos de Construcción y versión corregida de Transición.\\\hline
E6 & Plan y registros de pruebas, problemas, revisiones y evaluación. Equipo y director. & A lo largo del desarrollo y consolidación en Transición.\\\hline
E7 & Manual de instalación y uso, guía técnica e informe final con demostración. Director y usuarios previstos. & Preparación y cierre de Transición.\\\hline
\end{longtable}

Los documentos que permiten verificar la aplicación de estándares forman parte de los entregables y no se reducen a una mención bibliográfica. La tabla siguiente identifica su ubicación lógica, mientras que la sección \ref{sec:spmp-12-3} establece cómo evolucionarán sus versiones. Una evidencia podrá respaldar varios estándares mediante referencias, evitando copias independientes con resultados contradictorios. La selección representa una adaptación académica y no una certificación de conformidad~\cite{std29148,std12207,std29119p2,std29119p3,std29119p4,uml,ieee730}.
\begin{longtable}{|p{0.1\dimexpr\linewidth-34pt\relax}|p{0.67\dimexpr\linewidth-34pt\relax}|p{0.23\dimexpr\linewidth-34pt\relax}|}
\caption{Documentos para evaluar la aplicación de estándares. Fuente: adaptación de estándares del proyecto~\cite{adaptacion2026}.}\\
\hline
\textbf{ID} & \textbf{Documento y relación con entregables} & \textbf{Base}\\\hline
\endfirsthead
\hline
\textbf{ID} & \textbf{Documento y relación con entregables} & \textbf{Base}\\\hline
\endhead
D01 & Especificación, matriz de trazabilidad y registro de cambios de E1, consolidados hasta la entrega final. & 29148.\\\hline
D02 & Matriz de procesos, responsables y evidencias que vincula este PMP con la ejecución de E1--E7. & 12207.\\\hline
D03 & Modelos versionados de E2 y lista de comprobación de su coherencia con requisitos e implementación. & UML 2.5.1.\\\hline
D04 & Plan de pruebas, matriz requisito--caso y registros de ejecución e incidencias de E6. & 29119, partes 2, 3 y 4.\\\hline
D05 & Plan de calidad integrado en la sección \ref{sec:spmp-12-5} e informe de seguimiento por versiones de E6. & IEEE 730.\\\hline
\end{longtable}
